\scanpage{30}

\mspage{14}

les \unsure{solides|soides|valeurs} généraux, s'identifie au revêtement des orientations locales du précédent). Quand on se donne une jauge $e_M\in(\det M)^*$, on peut aussi considérer l'ens des repères symplectiques

\[
\underset{\substack{\shortparallel\\\mathrm{Rep.sp.}_\Lambda(M,e_M)/\{\pm\mathrm{id}_M\}}}{\operatorname{Repsp}_\Lambda(M,e_M)'}\subset\operatorname{Rep}_\Lambda(M)'\tag*{(75)}
\]

qui est un torseur sous $\mathrm{SL}(2,\Lambda)'$, et qui devient donc un ens. sur lequel $\mathrm{SL}(2,\mathbb{Z})'$ opère. Il lui correspond donc une structure $2$-polyédrale orientée, qu'il faudrait encore expliciter en termes des constructions \unsure{précédantes|précédentes} \deleted{\unsure{d}}(du \S\ I). Sauf erreur, pour $\Lambda$ de la forme $\mathbb{Z}$ ou $\mathbb{Z}/n\mathbb{Z}$ (ce qui est nécessaire pour que lesdites constructions fournissent bien un polyèdre, au sens strict \unsure{du terme|des termes}), on retrouve,

% archive sheet number: 29
